\section{Related Work}
\label{sec:related}

We organize the related work into three categories: studies on nested scheduling in ROS~2, efforts to deploy ROS~2 on real-time platforms, and the real-time capabilities of the RTEMS operating system.

\subsection{Nested Scheduling in ROS~2}

The scheduling semantics of the default ROS~2 Executor have attracted significant research attention due to the priority inversion and nondeterminism that arise from its nested scheduling architecture.

Casini et al.~\cite{casini2019} were the first to identify the priority inversion problem inherent in the ROS~2 Executor, demonstrating that the single-threaded Executor's FIFO-based dispatch of ready callbacks can cause lower-priority callbacks to block higher-priority ones. Tang et al.~\cite{tang2020} developed response time analysis for processing chains in ROS~2, providing analytical bounds for end-to-end latencies under the default Executor semantics. Choi et al.~\cite{choi2021} proposed PiCAS, a priority-driven chain-aware scheduling policy that reorders callback dispatch within the Executor to respect chain-level priorities, partially mitigating the priority inversion. Blass et al.~\cite{blass2021} extended the analysis by deriving response time bounds under starvation-freedom constraints, clarifying the trade-off between fairness and timeliness in the Executor. Jiang et al.~\cite{jiang2022} investigated real-time scheduling on multithreaded Executors, analyzing how the interplay between the Executor's internal dispatch logic and the OS scheduler affects schedulability.

More recently, Ishikawa-Aso et al.~\cite{ishikawa2025} proposed the CallbackIsolatedExecutor, which establishes a one-to-one mapping between callbacks and Linux pthreads, thereby delegating scheduling decisions entirely to the OS scheduler. While this approach eliminates the nested scheduling layer, it relies on the Linux CFS/SCHED\_FIFO scheduler, which does not provide deterministic interrupt latency and introduces significant per-callback overhead due to user--kernel mode transitions. Teper et al.~\cite{teper2024} presented the EventsExecutor, which bridges ROS~2 callback dispatch with classic real-time scheduling theory by exposing callback events as schedulable entities. Sun et al.~\cite{sun2024} revealed the root causes of nondeterminism in the ROS~2 Executor through systematic measurement, identifying the interplay between the wait-set mechanism and the Executor dispatch loop as the primary source of timing variability. Voronov et al.~\cite{voronov2024} uncovered timing anomalies in the rclcpp Executor, showing that certain execution orders can lead to response times far worse than predicted by conventional analysis.

Our work differs from the above approaches in a fundamental way: rather than attempting to fix or work around the nested scheduling problem within the existing Linux-based ROS~2 framework, we eliminate it entirely by porting ROS~2 to an open-source RTOS (RTEMS) and redesigning the Executor so that each callback is dispatched as an independent RTEMS task under the direct control of a deterministic fixed-priority preemptive scheduler.

\subsection{ROS~2 on Real-Time Platforms}

Several efforts have explored running ROS~2 on platforms other than standard Linux, but each exhibits significant limitations for real-time deployments.

The micro-ROS framework~\cite{micrROS} adapts ROS~2 for microcontrollers by using the rclc executor on NuttX and FreeRTOS. However, micro-ROS supports only a feature subset of the full ROS~2 stack and does not provide the rclcpp C++ API, which precludes the use of the vast majority of existing ROS~2 packages and necessitates rewriting application code against the C-based rcl API. Commercial RTOSes such as QNX and VxWorks offer ROS~2 ports with real-time guarantees, but their closed-source nature severely limits research reproducibility and community-driven improvement. The PREEMPT\_RT patch set for Linux provides improved real-time behavior, yet the underlying Linux kernel still exhibits non-deterministic interrupt latency, and the Completely Fair Scheduler (CFS) can interfere with real-time task execution~\cite{reghenzani2019}. OpenVela~\cite{openvela} supports multiple RTOS kernels but provides only limited ROS~2 integration, lacking a complete rclcpp port or executor redesign.

In contrast, our work represents the first full rclcpp-based ROS~2 port to an open-source RTOS (RTEMS), accompanied by a complete executor redesign that replaces the nested scheduling architecture with OS-native callback dispatch. This combination preserves full ROS~2 API compatibility while achieving deterministic real-time behavior.

\subsection{RTEMS Real-Time Capabilities}

RTEMS (Real-Time Executive for Multiprocessor Systems)~\cite{rtems} is an open-source real-time operating system that provides several capabilities critical for our approach.

The RTEMS default scheduler implements fixed-priority preemptive scheduling, which directly corresponds to the classic rate-monotonic scheduling model used in real-time analysis. RTEMS supports the priority inheritance protocol through the \texttt{RTEMS\_INHERIT\_PRIORITY} attribute, enabling automatic priority propagation to avoid unbounded priority inversion when callbacks share resources. The RTEMS watchdog facility employs a red-black tree timer scheduling mechanism that provides $O(\log n)$ insertion and removal of timers with bounded and deterministic overhead, a critical property for the precise periodic dispatch of timer callbacks. RTEMS also offers broad POSIX compliance, which significantly facilitates application portability and reduces the engineering effort required to port complex middleware such as ROS~2.

RTEMS has been extensively deployed in safety-critical domains including aerospace, defense, and medical systems, and provides certification artifacts for standards such as DO-178C and IEC~61508. This established certification track record makes our approach directly applicable to systems with stringent safety requirements, a property that neither Linux-based solutions nor closed-source alternatives can readily offer.
